\section{Selecting the Standard-Model orbit}
\label{sec:selection}

The Standard-Model line $\Phi_0$ lies in the $(\mathbf{10},\mathbf1,\mathbf3)$ component, which is small enough to be analysed completely. We first classify
the renormalizable potential on that component and give sufficient and necessary conditions for $\Phi_0$ to be its global minimum; then we add the gauge
and matter fluctuations;
and finally we return to the whole of $\Kc$.

\subsection{The tensor model}

Let $\su4\oplus\su2_R$ act on $\C^8=\C^4\otimes\C^2$ by Hermitian generators normalised by $\Tr T_aT_b=\tfrac12\delta_{ab}$ in the fundamental, and on complex
symmetric $8\times8$ matrices by $S\mapsto USU^{\!\top}$. The subspace
\[
\mathcal T=\mathrm{Sym}^2(\C^4)\otimes\mathrm{Sym}^2(\C^2)\subset\mathrm{Sym}^2(\C^8),\qquad \dim_\C\mathcal T=30,\qquad |S|^2=\Tr S^\dagger S,
\]
is invariant and irreducible, isomorphic to $(\mathbf{10},\mathbf1,\mathbf3)$. With $e_1,\dots,e_4$ and $e_\pm$ the standard bases of $\C^4$ and $\C^2$, $\BmL=\operatorname{diag}(\tfrac13,\tfrac13,\tfrac13,-1)$
and $T^R_3=\operatorname{diag}(\tfrac12,-\tfrac12)$, the vectors $e_i\otimes e_+$, $e_4\otimes e_+$, $e_i\otimes e_-$, $e_4\otimes e_-$ ($i\le3$) carry $Y=\tfrac23,0,-\tfrac13,-1$:
they are $u_R$, $\nu_R$, $d_R$, $e_R$, and a component of $S$ is a pair of them. Put $u_0=e_4\otimes e_+$ and $S_0=u_0u_0^{\!\top}$.

\begin{proposition}[From kernels to tensors]
\label{prop:tensor}
The component $\Kc_{(\mathbf{10},\mathbf1,\mathbf3)}$ is $J$-invariant of real dimension $60$, and with $J$ as complex structure it is isomorphic to $\mathcal T$ by a
Pati--Salam-equivariant $\iota$ with $\Tr\Phi^{\!\top}\Phi=4|\iota(\Phi)|^2$, unique up to a phase, with $\iota(\Phi_0)=e^{i\varphi}S_0$. For
$S=\iota(\Phi)$, $\Phi^{\!\top}\Phi$ is supported on the $\Gamma=+1$ half and has the spectrum of $S^\dagger S$, each eigenvalue four times, so
\[
N:=\Tr\Phi^{\!\top}\Phi=4|S|^2,\qquad Q:=\Tr(\Phi^{\!\top}\Phi)^2=4\Tr(S^\dagger S)^2,\qquad \frac Q{N^2}=\frac14\,\frac{\Tr(S^\dagger S)^2}{|S|^4}.
\]
\end{proposition}

Let $P_1,\dots,P_5$ be the orthogonal projectors of $\mathrm{Sym}^2\mathcal T$ onto its five irreducible components
\[
(\mathbf{35},\mathbf5)\oplus(\mathbf{35},\mathbf1)\oplus(\mathbf{20}',\mathbf5)\oplus(\mathbf{20}',\mathbf1)\oplus(\mathbf{45},\mathbf3),\qquad 175+35+100+20+135=465,
\]
and put $I_k(S)=\langle S\otimes S,P_k(S\otimes S)\rangle\ge0$, so that $\sum_kI_k=|S|^4$.

\begin{theorem}[The renormalizable potential of $(\mathbf{10},\mathbf1,\mathbf3)$]
\label{thm:potential}
The most general renormalizable invariant potential on $\mathcal T$ is
\begin{equation}
\label{eq:Vgeneral}
V(S)=-\mu^2|S|^2+\sum_{k=1}^5\lambda_kI_k(S)+\kappa\,\mathcal D(S)+\bar\kappa\,\overline{\mathcal D(S)},\qquad \lambda_k\in\R,\ \kappa\in\C,
\end{equation}
where the holomorphic quartic is
\[
\mathcal D(S)=\tfrac1{24}\,\epsilon^{abcd}\epsilon^{efgh}P_{ab,ef}(S)P_{cd,gh}(S),\qquad P_{ab,ef}(S)=\varepsilon_{su}\varepsilon_{tv}S_{as,et}S_{bu,fv},
\]
with colour indices $a,\dots,h$ and $\su2_R$ indices $s,t,u,v$. There is no invariant cubic. Moreover:
\begin{enumerate}[label=(\roman*),leftmargin=2.2em]
\item $\mathcal D(e^{i\theta}S)=e^{4i\theta}\mathcal D(S)$, every $I_k$ is phase invariant, and $\mathcal D$ is not a combination of the $I_k$.
\item $\mathcal D$ depends on $S$ only through the $(\mathbf{20}',\mathbf1)$ component of $S\otimes S$, and $|\mathcal D(S)|\le\tfrac38I_4(S)$, with equality at
$S=\id_4\otimes\id_2$.
\item Under fermion number $S\mapsto e^{2i\theta}S$, so the $I_k$ are invariant and $\mathcal D\mapsto e^{8i\theta}\mathcal D$: \eqref{eq:Vgeneral} is invariant
under fermion number if and only if $\kappa=0$. $\mathcal D$ has $\BmL=0$ and fermion number $8$, hence baryon number $2$.
\end{enumerate}
\end{theorem}

The trivial representation occurs in $\mathcal T\otimes\bar{\mathcal T}$, $\mathrm{Sym}^2\mathcal T$, $\mathrm{Sym}^3\mathcal T$, $\mathrm{Sym}^2\mathcal T\otimes\bar{\mathcal T}$,
$\mathrm{Sym}^4\mathcal T$, $\mathrm{Sym}^3\mathcal T\otimes\bar{\mathcal T}$ and $\mathrm{Sym}^2\mathcal T\otimes\mathrm{Sym}^2\bar{\mathcal T}$ with multiplicities $1,0,0,0,1,0,5$,
which is the content of the theorem. The phase of $\kappa$ can be removed by $S\mapsto e^{i\alpha}S$, so only $|\kappa|$ matters. The contraction of four
$(\mathbf{10},\mathbf1,\mathbf3)$ fields with two $SU(4)$ $\epsilon$ tensors is the baryon-number-violating coupling of Pati--Salam models through which diquark
components mediate neutron--antineutron oscillation \cite{MohapatraMarshak80,ChackoMohapatra99}.

Two lemmas organise the quartics. The first is the recoupling identity $\Tr(S^\dagger S)^2=I_1-\tfrac12I_2-\tfrac12I_3+\tfrac14I_4+0\cdot I_5$, obtained from
Schur--Weyl duality on $(\C^4)^{\otimes4}\otimes(\C^2)^{\otimes4}$. The second characterises the orbit of $S_0$.

\begin{lemma}[The highest-weight orbit]
\label{lem:hw}
For $S\in\mathcal T\setminus0$ the following are equivalent: $S\otimes S$ lies in $(\mathbf{35},\mathbf5)$; $\Tr(S^\dagger S)^2=|S|^4$; $\operatorname{rank}S=1$;
$S=(a\otimes b)(a\otimes b)^{\!\top}$ with $a\in\C^4$, $b\in\C^2$; $S$ lies on the orbit $\mathcal O_{\rm hw}$ of $\C^\times S_0$.
\end{lemma}

This is the Kostant--Lichtenstein description of the highest-weight orbit by quadrics \cite{Lichtenstein82}. On $\mathcal O_{\rm hw}$ the colour factor of $S$
has rank one, so $\mathcal D$ vanishes there; more precisely $\mathcal D$ and its first and second derivatives vanish at $S_0$, and
\[
\mathcal D(s_0S_0+\delta)=s_0\Bigl(\tfrac13\epsilon_{ijk}\epsilon_{lmn}U_{il}D_{jm}D_{kn}-\tfrac1{12}\epsilon_{ijk}\epsilon_{lmn}M_{il}M_{jm}D_{kn}\Bigr)+\mathcal D(\delta),
\]
with $U$, $M$, $D$ the symmetric $3\times3$ matrices of the colour sextets $u_Ru_R$, $u_Rd_R$, $d_Rd_R$. In the vacuum the holomorphic quartic contributes
nothing to the masses and produces the cubic couplings $\kappa v\,\epsilon\epsilon\,\Delta_{uu}\Delta_{dd}\Delta_{dd}$ and $\kappa v\,\epsilon\epsilon\,\Delta_{ud}\Delta_{ud}\Delta_{dd}$ of
diquarks of hypercharges $\tfrac43,\tfrac13,-\tfrac23$, each of baryon number two. With the Yukawa couplings of the diquarks to quark pairs, here the
single coupling $y$ of \eqref{eq:Saction}, the vertex closes into a six-quark $\Delta B=2$ operator, the Pati--Salam mechanism of neutron--antineutron
oscillation; the present limits on the oscillation time are $8.6\times10^7$\,s for free neutrons \cite{BaldoCeolin94} and $4.7\times10^8$\,s from
neutrons bound in oxygen \cite{SuperK21}.

\subsection{Where the Standard-Model orbit is the global minimum}

Number the components $k=1,\dots,5$ as above, put $f_k(S)=I_k(S)/|S|^4$ (so $f_k\ge0$, $\sum f_k=1$), $d(S)=|\mathcal D(S)|/|S|^4$, and
\[
\Delta_k=\lambda_k-\lambda_1\ (k\ge2),\qquad \Lambda(S)=\sum_k\lambda_kf_k(S)=\lambda_1+\sum_{k\ge2}\Delta_kf_k(S),\qquad \Lambda_\kappa(S)=\Lambda(S)-2|\kappa|\,d(S).
\]
For a unit vector $\hat S$, $\min_\varphi V(re^{i\varphi}\hat S)=-\mu^2r^2+\Lambda_\kappa(\hat S)r^4$, because the $I_k$ do not depend on the phase and
$2\Re(\kappa e^{4i\varphi}\mathcal D(\hat S))$ reaches $-2|\kappa||\mathcal D(\hat S)|$. With $\mu^2>0$ the potential is bounded below iff $\Lambda_\kappa>0$ on the unit
sphere, the minimum along a ray is $-\mu^4/(4\Lambda_\kappa)$, and the vacuum is the configuration of smallest $\Lambda_\kappa$.

\begin{theorem}[Selection without the holomorphic quartic]
\label{thm:select0}
Let $\kappa=0$, $\mu^2>0$, $\lambda_1>0$ and $\Delta_2,\dots,\Delta_5>0$. The global minima of \eqref{eq:Vgeneral} are exactly the points of $\mathcal O_{\rm hw}$ with
$|S|^2=\mu^2/(2\lambda_1)$, with $V_{\min}=-\mu^4/(4\lambda_1)$. They form one orbit of real dimension nine, and the stabiliser of each point is conjugate to
$\su3\oplus\uu1_Y$ inside $\su4\oplus\su2_R$ with $Y=T^R_3+\tfrac12(\BmL)$; together with $\su2_L$, which acts trivially on $(\mathbf{10},\mathbf1,\mathbf3)$, this is the
Standard-Model algebra.
\end{theorem}

\begin{proof}
$\Lambda\ge\lambda_1$, with equality iff $f_2=\dots=f_5=0$, that is iff $S\in\mathcal O_{\rm hw}$ by \cref{lem:hw}; the unit-norm part of $\mathcal O_{\rm hw}$ is a single
orbit, of dimension $18-9$.
\end{proof}

The unbroken $\uu1$ is an output, computed as the stabiliser: the weights of the vacuum, not a choice of embedding, fix the hypercharge. This is the
vacuum of Pati--Salam and left--right models with a $(\mathbf{10},\mathbf1,\mathbf3)$ field \cite{PatiSalam,MohapatraSenjanovic81}; the theorem says for which couplings
the dynamics chooses it.

\begin{theorem}[Selection with the holomorphic quartic]
\label{thm:selectk}
Let $\mu^2>0$ and $\kappa\in\C$.
\begin{enumerate}[label=(\alph*),leftmargin=2.2em]
\item If $\lambda_1,\lambda_2,\lambda_3,\lambda_5>0$ and $\lambda_4>\tfrac34|\kappa|$, the potential is bounded below; if it is bounded below then $\lambda_1>0$ and
$2\lambda_1+\lambda_2+2\lambda_3+\lambda_4>\tfrac34|\kappa|$.
\item If $\lambda_1>0$, $\Delta_2,\Delta_3,\Delta_5>0$ and $\Delta_4>\tfrac34|\kappa|$, the global minima are exactly the points of $\mathcal O_{\rm hw}$ with
$|S|^2=\mu^2/(2\lambda_1)$.
\item If those points are global minima, then $\Delta_2\ge0$, $\Delta_3\ge0$, $2\Delta_2+2\Delta_3+4\Delta_4+9\Delta_5\ge0$ and
$\Delta_2+2\Delta_3+\Delta_4\ge\tfrac34|\kappa|$.
\item The second variation at the Standard-Model orbit does not depend on $\kappa$.
\end{enumerate}
\end{theorem}

At the benchmark couplings $\lambda=(1,1.7,2.3,3.1,2.9)$ the bounds read: $\tfrac43\Delta_4=2.8$, sufficient for selection; $\tfrac43(\Delta_2+2\Delta_3+\Delta_4)=7.2$,
necessary for selection and equal to the computed threshold; $\tfrac43\lambda_4=4.13$, sufficient for boundedness; and $\tfrac43(2\lambda_1+\lambda_2+2\lambda_3+\lambda_4)=15.2$,
necessary for boundedness. Gradient flows confirm the transition at $|\kappa|=7.2$: from random starts the flow reaches the Standard-Model orbit for every
$|\kappa|$, as it must, since the orbit is always a local minimum, but for $6\le|\kappa|\le7.5$ one of two starts ends on the orbit of $\id_4\otimes\id_2$, which is
the higher of the two minima for $|\kappa|\le7$ and the lower for $|\kappa|\ge7.5$: at the benchmark,
$f(\id_4\otimes\id_2)=(\tfrac13,\tfrac16,\tfrac13,\tfrac16,0)$ and $d(\id_4\otimes\id_2)=\tfrac1{16}$, so
$\Lambda_\kappa(\id_4\otimes\id_2)-\Lambda_\kappa(S_0)=0.9-\tfrac18|\kappa|$, which changes sign at $|\kappa|=7.2$.

\begin{theorem}[Second variation at the Standard-Model orbit]
\label{thm:secondvar}
At $S_v=vS_0$ the sixty real fluctuations split into $SU(3)\times U(1)_Y$ multiplets, each with one mass:
\begin{center}
\small
\begin{tabular}{@{}llccl@{}}
\toprule
pair & $SU(3)_Y$ & $\dim_\R$ & $m^2/v^2$ & role\\
\midrule
$\nu_R\nu_R$ (modulus) & $1_0$ & $1$ & $4\lambda_1$ & radial\\
$\nu_R\nu_R$ (phase), $\nu_Re_R$, $u_R\nu_R$ & $1_0$, $1_{-1}$, $3_{2/3}$ & $1+2+6$ & $0$ & eaten\\
$e_Re_R$ & $1_{-2}$ & $2$ & $\tfrac43\Delta_2$ & $\Delta_{ee}$\\
$u_Re_R+d_R\nu_R$ & $3_{-1/3}$ & $6$ & $\Delta_5$ & $\Delta_{ue}$\\
$d_Re_R$ & $3_{-4/3}$ & $6$ & $\tfrac23\Delta_2+\Delta_5$ & $\Delta_{de}$\\
$u_Ru_R$ & $6_{4/3}$ & $12$ & $\tfrac43\Delta_3$ & $\Delta_{uu}$\\
$u_Rd_R$ & $6_{1/3}$ & $12$ & $\tfrac23\Delta_3+\Delta_5$ & $\Delta_{ud}$\\
$d_Rd_R$ & $6_{-2/3}$ & $12$ & $\tfrac49\Delta_2+\tfrac29\Delta_3+\tfrac49\Delta_4+\Delta_5$ & $\Delta_{dd}$\\
\bottomrule
\end{tabular}
\end{center}
Consequently $S_v$ is a non-degenerate local minimum on $\mathcal T$ iff $\lambda_1,\Delta_2,\Delta_3,\Delta_5>0$ and $2\Delta_2+2\Delta_3+4\Delta_4+9\Delta_5>0$, a region that
strictly contains the cone of \cref{thm:select0}; and, independently of the couplings,
\begin{equation}
\label{eq:sumrules}
m_{de}^2=m_{ue}^2+\tfrac12m_{ee}^2,\qquad m_{ud}^2=m_{ue}^2+\tfrac12m_{uu}^2 .
\end{equation}
\end{theorem}

\begin{example}[Checking a sum rule]
\label{ex:sumrule}
At the benchmark $\lambda=(1,1.7,2.3,3.1,2.9)$ one has $\Delta=(0.7,1.3,2.1,1.9)$, so in units of $v^2$: $m_{ee}^2=\tfrac43\cdot0.7=0.9\overline3$,
$m_{ue}^2=1.9$, $m_{de}^2=\tfrac23\cdot0.7+1.9=2.3\overline6$, and indeed $1.9+\tfrac12\cdot0.9\overline3=2.3\overline6$. Likewise $m_{uu}^2=\tfrac43\cdot1.3=1.7\overline3$
and $m_{ud}^2=\tfrac23\cdot1.3+1.9=2.7\overline6=1.9+\tfrac12\cdot1.7\overline3$. The full $60\times60$ mass matrix at these couplings and $\mu^2=2$ has the
predicted spectrum to $4\times10^{-15}$.
\end{example}

The colour-sextet components $\Delta_{uu},\Delta_{ud},\Delta_{dd}$ are the diquarks just met, and $\Delta_{ee}$ is the doubly charged scalar of left--right models.
The sum rules \eqref{eq:sumrules} hold for the most general renormalizable potential of a single $(\mathbf{10},\mathbf1,\mathbf3)$, with or without $\kappa$. When one of the
bounds $\Delta_2,\Delta_3,\Delta_5>0$ fails, gradient flows end on configurations of lower energy that break electric charge or colour: the configuration
$(e_4e_4)(e_+e_-)$, for instance, keeps colour, but its unbroken abelian factor is $T^R_3$ instead of $Y$.

\subsection{The trace potential}

The kernel field carries two quartic invariants that do not refer to the Pati--Salam decomposition at all, $N^2$ and $Q=\Tr(\Phi^{\!\top}\Phi)^2$; they are
invariant under every orthogonal change of basis of the block that preserves $\Kc$, in particular under fermion number, so the trace potential has
$\kappa=0$.

\begin{theorem}[The trace potential on $(\mathbf{10},\mathbf1,\mathbf3)$]
\label{thm:tracepot}
Let $V_{\rm tr}=-\mu^2N+\alpha N^2+\beta Q$ with $\mu^2>0$. In the variables of \cref{thm:potential}, $\lambda_k=16\alpha+4\beta a_k$ with $a=(1,-\tfrac12,-\tfrac12,\tfrac14,0)$,
so $(\Delta_2,\Delta_3,\Delta_4,\Delta_5)=-4\beta\,(\tfrac32,\tfrac32,\tfrac34,1)$. If $\beta<0$ and $\alpha+\tfrac14\beta>0$, the global minimum is the Standard-Model orbit; if
$\beta>0$ and $\alpha+\tfrac1{32}\beta>0$, the minima are the $S$ with $S^\dagger S=\tfrac18|S|^2\id_8$, and every one of them breaks colour; if $\beta=0$ the potential is
$O(60)$-symmetric and the orbit is chosen by fluctuations.
\end{theorem}

$Q/N^2$ measures how concentrated the Majorana mass matrix is: it equals $\tfrac14$ when one fermion carries all of it and $\tfrac1{32}$ when the eight
right-handed fermions share it equally. A negative $\beta$ rewards concentration on a single line, and the only line on which the dynamics can concentrate
without breaking colour is $\nu_R$. The Standard-Model vacuum is the concentrated phase of the Majorana mass, which is the three-constituent form of
the Majorana phase of \cref{thm:twophases}.

\subsection{The fluctuations decide: back-reaction on \texorpdfstring{$(\mathbf{10},\mathbf1,\mathbf3)$}{(10,1,3)}}

Take the tree potential degenerate on $(\mathbf{10},\mathbf1,\mathbf3)$, $\lambda_1=\dots=\lambda_5$ and $\kappa=0$, the case $\beta=0$. Every unit direction of $\mathcal T$ is then a
tree-level minimum, and the orbit is chosen by the one-loop potential \eqref{eq:V1}. Gauge and Yukawa interactions respect fermion number, so they do not
generate $\kappa$ in perturbation theory. Write $L=\ln(\Lambda_{\rm UV}^2/v^2)$. The orbit-dependent part of the leading logarithm is
\[
\delta V=\frac L{64\pi^2}\Bigl[2\sum_{\rm fam}\Tr(M_F^\dagger M_F)^2-3\Tr M_V^4-\Tr M_S^4\Bigr],
\]
where for the degenerate tree potential $\Tr M_S^4$ depends on $|S|$ only. Here and in \cref{sec:fullK} the field is normalised canonically. By
\cref{prop:tensor}, $\Tr\Phi^{\!\top}\Phi=4|\iota(\Phi)|^2$, so the kinetic term of \cref{ass:field} is $\Tr(DS)^\dagger(DS)$ for $S:=\sqrt2\,\iota(\Phi)$, and from
here on $S$ denotes this field and $v=|S|=|\Phi|/\sqrt2=\sqrt2\,r$, the $v$ of \cref{sec:fullK,sec:matter}. In these variables the right-handed neutrino mass
$|y|r$ of \cref{sec:kernels} reads $(|y|/\sqrt2)\,v$, so the $y$ of $M_F=yS$ below is the coupling of \eqref{eq:Saction} divided by $\sqrt2$, while the
coupling $g$ is the same as in \cref{thm:compat}. The vector mass matrix is $(M_V^2)_{ab}=2g^2\Re\Tr(T_aS+ST_a^{\!\top})^\dagger(T_bS+ST_b^{\!\top})$
with one coupling for $\su4$ and $\su2_R$, and the eight right-handed Weyl fermions get the Majorana matrix $M_F=yS$.

\begin{theorem}[One-loop selection on $(\mathbf{10},\mathbf1,\mathbf3)$]
\label{thm:oneloop1013}
At leading logarithm, $\Tr M_V^4=g^4\sum_kc_kI_k$ with $c=(33,3,-3,39,5)$ and $\Tr(M_F^\dagger M_F)^2=y^4\sum_ka_kI_k$, and the fluctuations generate the gaps
\[
\Delta_k=\frac L{64\pi^2}\,g^4\bigl(G_k-x\,F_k\bigr),\qquad G=(90,108,-18,84),\quad F=(3,3,\tfrac32,2),\qquad x=\frac{\sum_{\rm fam}y^4}{g^4}.
\]
Gauge fluctuations favour the Standard-Model orbit in the directions $k=2,3,5$, and matter fluctuations disfavour it in every direction. The Standard-Model
orbit is a non-degenerate local minimum on $\mathcal T$ iff $x<30$, the binding conditions being those of $\Delta_{ee}$ and $\Delta_{dd}$. The configuration $\id_4\otimes\id_2$
has excess $\tfrac L{64\pi^2}g^4(48-\tfrac74x)$ over the Standard-Model orbit, negative for $x>\tfrac{192}7\approx27.43$.
\end{theorem}

At $x=0$ the radiatively generated masses are, in units of $\tfrac L{64\pi^2}g^4v^2$, $m_{ee}^2=120$, $m_{ue}^2=84$, $m_{de}^2=144$, $m_{uu}^2=144$, $m_{ud}^2=156$ and
$m_{dd}^2=120$, which satisfy \eqref{eq:sumrules}; the fifty massive scalars are pseudo-Goldstone bosons of the accidental $O(60)$ symmetry of the tree potential
\cite{Weinberg72,GildenerWeinberg76}. A projected gradient flow on the unit sphere of $\mathcal T$ from random starts finds no configuration below the Standard-Model
orbit for $x\le27$, finds the orbit of $\id_4\otimes\id_2$ as the only competitor, and confirms that its energy crosses zero at $x=192/7$; between $192/7$ and $30$ the
Standard-Model vacuum is metastable on $\mathcal T$. We conjecture that on $\mathcal T$ the Standard-Model orbit is the global minimum of the one-loop potential exactly
for $x<192/7$.

With $M_N=yv$ and $M_{W_R}=gv$, $x=\sum_{\rm fam}(M_N/M_{W_R})^4$; in the couplings of \eqref{eq:Saction} and \cref{thm:compat},
$x=\sum_{\rm fam}y^4/(4g^4)$. Stated in terms of masses, $x$ does not depend on the normalisation. The content of the bound is structural: the Majorana Yukawa coupling, matter's only handle on the
vacuum, pushes towards configurations in which the Majorana mass is shared among several fermions, and the gauge fluctuations hold the Standard-Model orbit only
while the heavy neutrinos are not much heavier than $W_R$.

\subsection{The full kernel space}
\label{sec:fullK}

The general invariant potential on all of $\Kc$ has too many couplings to be classified usefully: the numbers of linearly independent Pati--Salam-invariant
real polynomials of degrees $1,2,3,4$ are $0,18,0,686$, and with fermion number imposed $0,12,0,402$. Two structures do not depend on these couplings.

\begin{proposition}[A left--right involution]
\label{prop:LR}
Let $O=E_1E_7E_8$ restricted to $V_*$. It commutes with $P_*$ and with the rotations, satisfies $O^{\!\top}K^0O=K^0$, $O^{\!\top}K^iO=-K^i$, $O\Gamma O^{\!\top}=-\Gamma$ and
$OJO^{\!\top}=J$, and maps $t^L_i$ to $\pm t^R_i$ and each colour bivector to $\pm$ itself. Hence $\psi(t,\mathbf x)\mapsto O\psi(t,-\mathbf x)$, $\Phi\mapsto O\Phi O^{\!\top}$,
with the corresponding map of the gauge fields, is a symmetry of the kinetic, Yukawa and gauge terms of \eqref{eq:Saction}, and of the trace invariants. The image
$\Phi_0'=O\Phi_0O^{\!\top}$ is a kernel in $(\mathbf{10},\mathbf3,\mathbf1)$ orthogonal to $\Phi_0$, the mirror vacuum of left--right models, and along the great circle
$\hat u(\theta)=\cos\theta\,\hat u_0+\sin\theta\,\hat u_0'$, $Q=\tfrac14-\tfrac18\sin^22\theta$.
\end{proposition}

The Standard-Model vacuum and its mirror are exchanged by a symmetry of the action, so an $O$-invariant potential cannot prefer one of them; the choice between
breaking $\su2_R$ and breaking $\su2_L$ is spontaneous, as in models of spontaneous parity violation \cite{SenjanovicMohapatra75}.

Now take the tree potential degenerate on all of $\Kc$, $V_0=-\mu^2N+\alpha N^2$, invariant under the accidental $O(272)$. For a unit vector $\hat u\in\Kc$ put
\[
\begin{gathered}
Q(\hat u)=\Tr(\hat u^{\!\top}\hat u)^2,\qquad T(\hat u)=\Tr\bigl(G^{-1}K(\hat u)\bigr)^2,\qquad K(\hat u)_{ab}=\langle[t_a,\hat u],[t_b,\hat u]\rangle,\\
E_x(\hat u)=8x\,Q(\hat u)-\tfrac3{16}T(\hat u),
\end{gathered}
\]
with the $21$ Pati--Salam generators $t_a$ and the module metric $G$. The orientation-dependent part of the leading logarithm is then
$\delta V=\tfrac L{64\pi^2}g^4v^4E_x(\Phi/|\Phi|)$ with $v=|\Phi|/\sqrt2$, and at the Standard-Model point $E_x(\hat u_0)=2x-99$.

\begin{theorem}[One-loop second variation on all of $\Kc$]
\label{thm:fullsecond}
Let $H(x)$ be the second variation of $E_x$ at $\hat u_0$ on the $271$-dimensional tangent space of the unit sphere.
\begin{enumerate}[label=(\roman*),leftmargin=2.2em]
\item $H(x)=H_G+xH_F$ with $H_F=-8(\id-\Pi_\nu)$, where $\Pi_\nu$ is the orthogonal projector of rank $32$ onto the kernels that pair $\nu_R$ with a Weyl fermion of the
block. $H_G$ preserves each Pati--Salam component of $\Kc$.
\item With $m^2=\tfrac12\times$(eigenvalue of $H(x)$) in units of $\tfrac L{64\pi^2}g^4v^2$, the $271$ angular modes have the masses of \cref{tab:271}.
\item $H(x)$ is non-negative with kernel exactly the nine Goldstone directions if and only if $x<27$. At $x=27$ the six modes of $\ell_L\ell_L$, the left-handed triplet,
become massless, and for $x>27$ they are tachyonic.
\end{enumerate}
\end{theorem}

\begin{table}[t]
\centering
\footnotesize
\setlength{\tabcolsep}{4pt}
\begin{tabular}{@{}llccl@{}}
\toprule
fermion pair & multiplet & $\dim_\R$ & component & $m^2$\\
\midrule
$\nu_R\nu_R$ (phase), $\nu_Re_R$, $\nu_Ru_R$ & $(1,1)_0,(1,1)_1,(3,1)_{2/3}$ & $1+2+6$ & $(\mathbf{10},\mathbf1,\mathbf3)$ & $0$ (Goldstone)\\
$e_Re_R$ & $(1,1)_2$ & $2$ & $(\mathbf{10},\mathbf1,\mathbf3)$ & $120-4x$\\
$d_Re_R$ & $(3,1)_{4/3}$ & $6$ & $(\mathbf{10},\mathbf1,\mathbf3)$ & $144-4x$\\
$u_Ru_R$ & $(6,1)_{4/3}$ & $12$ & $(\mathbf{10},\mathbf1,\mathbf3)$ & $144-4x$\\
$u_Rd_R$ & $(6,1)_{1/3}$ & $12$ & $(\mathbf{10},\mathbf1,\mathbf3)$ & $156-4x$\\
$d_Rd_R$ & $(6,1)_{2/3}$ & $12$ & $(\mathbf{10},\mathbf1,\mathbf3)$ & $120-4x$\\
$u_Re_R\pm d_R\nu_R$ & $(3,1)_{1/3}$ & $6+6$ & $(\mathbf{10},\mathbf1,\mathbf3)$, $(\mathbf6,\mathbf1,\mathbf1)_R$ & eigenvalues of $M_3$\\
$u_Rd_R$ (antisymmetric) & $(3,1)_{1/3}$ & $6$ & $(\mathbf6,\mathbf1,\mathbf1)_R$ & $180-4x$\\
$u_Le_L-d_L\nu_L$, $u_Ld_L$ & $(3,1)_{1/3}$ & $6+6$ & $(\mathbf6,\mathbf1,\mathbf1)_L$ & $180-4x$\\
$\ell_L\ell_L$ & $(1,3)_1$ & $6$ & $(\mathbf{10},\mathbf3,\mathbf1)$ & $108-4x$\\
$q_L\ell_L$ & $(3,3)_{1/3}$ & $18$ & $(\mathbf{10},\mathbf3,\mathbf1)$ & $156-4x$\\
$q_Lq_L$ & $(6,3)_{1/3}$ & $36$ & $(\mathbf{10},\mathbf3,\mathbf1)$ & $180-4x$\\
$q_L$--$u_R$, $q_L$--$d_R$ (octet) & $(8,2)_{1/2}$ & $64$ & $(\mathbf{15},\mathbf2,\mathbf2)$ & $174-4x$\\
$\ell_L$--$u_R$, $q_L$--$e_R$ & $(3,2)_{7/6}$ & $24$ & $(\mathbf{15},\mathbf2,\mathbf2)$ & $162-4x$\\
$\ell_L$--$d_R$ & $(3,2)_{1/6}$ & $12$ & $(\mathbf{15},\mathbf2,\mathbf2)$ & $114-4x$\\
$q_L$--$\nu_R$ & $(3,2)_{1/6}$ & $12$ & $(\mathbf{15},\mathbf2,\mathbf2)$ & $114$\\
$q_L$--$d_R$, $\ell_L$--$e_R$ (singlets) & $(1,2)_{1/2}$ & $4+4$ & $(\mathbf{15},\mathbf2,\mathbf2)$, $(\mathbf1,\mathbf2,\mathbf2)$ & $138-4x$, $186-4x$\\
$q_L$--$u_R$, $\ell_L$--$\nu_R$ (singlets) & $(1,2)_{1/2}$ & $4+4$ & $(\mathbf{15},\mathbf2,\mathbf2)$, $(\mathbf1,\mathbf2,\mathbf2)$ & eigenvalues of $M_2$\\
\bottomrule
\end{tabular}
\caption{The radiatively generated masses of the $271$ angular modes of the kernel field at the Standard-Model vacuum, in units of $\tfrac L{64\pi^2}g^4v^2$, with
$x=\sum_{\rm fam}(M_N/M_{W_R})^4$ (\cref{thm:fullsecond}). Pairs written $ff'$ are of Majorana type, pairs written $f$--$f'$ of Dirac type; multiplets are given
under $SU(3)\times SU(2)_L$ with $|Y|$. The two mixing matrices are
$M_3=\bigl(\begin{smallmatrix}84-2x&2x\\2x&180-2x\end{smallmatrix}\bigr)$ and $M_2=\bigl(\begin{smallmatrix}138-x&\sqrt3x\\\sqrt3x&186-3x\end{smallmatrix}\bigr)$, with eigenvalues
$132-2x\mp2\sqrt{x^2+576}$ and $162-2x\mp2\sqrt{x^2-12x+144}$.}
\label{tab:271}
\end{table}

Statement (i) is the mechanism by which matter opposes the Standard-Model orbit, now on all of $\Kc$. At fixed $|\Phi|$ the fermion loop $2\sum m_f^4$ is largest when the
whole Majorana mass sits on $\nu_R$; it decreases quadratically along every direction that moves mass onto pairs not containing $\nu_R$ and is stationary along the
directions that pair $\nu_R$ with another fermion. Every mode that does not involve $\nu_R$ therefore has slope $-4$, and the first of them to become massless is the one
with the smallest gauge contribution, the left-handed triplet. The mixing in $M_3$, between the $(\mathbf{10},\mathbf1,\mathbf3)$ mode $u_Re_R+d_R\nu_R$ and the $(\mathbf6,\mathbf1,\mathbf1)_R$
mode $u_Re_R-d_R\nu_R$, is an instance of the mixing announced at the end of \cref{sec:kernels}: it moves the $\Delta_{ue}$ threshold from $42$ to $315/11$. The
thresholds are the zeros of the lowest lines,
\[
x=27\ (\ell_L\ell_L),\qquad x=28.5\ (\ell_L\text{--}d_R),\qquad x=\tfrac{315}{11}\approx28.64\ (M_3),\qquad x=30\ (e_Re_R,\ d_Rd_R),
\]
and \cref{fig:radiative} shows them. Along the great circle of \cref{prop:LR},
\[
E_x(\hat u(\theta))=2x-99+(27-x)\sin^22\theta ,
\]
so for $x>27$ the Standard-Model vacuum descends towards the left--right symmetric configuration $\theta=\pi/4$, where $\nu_R$ and a left-handed lepton carry equal
Majorana masses.

\begin{figure}[t]
\centering
\begin{tikzpicture}
\begin{axis}[width=12.5cm,height=7.6cm,xmin=0,xmax=32,ymin=-15,ymax=200,
  xlabel={$x=\sum_{\rm fam}(M_N/M_{W_R})^4$},ylabel={$m^2$ in units of $\frac{L}{64\pi^2}g^4v^2$},
  axis lines=left,clip=true,xtick={0,5,10,15,20,25,30},ytick={0,50,100,150,200},
  ylabel style={font=\small},xlabel style={font=\small},
  legend style={at={(0.02,0.03)},anchor=south west,font=\footnotesize,fill=white,draw=mgray},legend cell align=left]
\fill[lgray!60] (axis cs:24.884,-15) rectangle (axis cs:27,200);
\addplot[gray!70,thin,forget plot,domain=0:32] {144-4*x};
\addplot[gray!70,thin,forget plot,domain=0:32] {156-4*x};
\addplot[gray!70,thin,forget plot,domain=0:32] {162-4*x};
\addplot[gray!70,thin,forget plot,domain=0:32] {174-4*x};
\addplot[gray!70,thin,forget plot,domain=0:32] {180-4*x};
\addplot[gray!70,thin,forget plot,domain=0:32] {186-4*x};
\addplot[gray!70,thin,forget plot,domain=0:32] {114};
\addplot[gray!70,thin,forget plot,domain=0:32,samples=60] {162-2*x+2*sqrt(x^2-12*x+144)};
\addplot[gray!70,thin,forget plot,domain=0:32,samples=60] {162-2*x-2*sqrt(x^2-12*x+144)};
\addplot[gray!70,thin,forget plot,domain=0:32,samples=60] {132-2*x+2*sqrt(x^2+576)};
\addplot[scol,very thick,domain=0:32] {108-4*x}; \addlegendentry{$\ell_L\ell_L$: $108-4x$}
\addplot[ocol,thick,domain=0:32] {114-4*x}; \addlegendentry{$\ell_L$--$d_R$: $114-4x$}
\addplot[gcol,thick,domain=0:32,samples=80] {132-2*x-2*sqrt(x^2+576)}; \addlegendentry{lower eigenvalue of $M_3$}
\addplot[tcol,thick,domain=0:32] {120-4*x}; \addlegendentry{$e_Re_R$, $d_Rd_R$: $120-4x$}
\draw[black,dashed] (axis cs:0,0) -- (axis cs:32,0);
\draw[gray,densely dotted] (axis cs:27,-15) -- (axis cs:27,200);
\draw[gray,densely dotted] (axis cs:28.5,-15) -- (axis cs:28.5,200);
\draw[gray,densely dotted] (axis cs:28.636,-15) -- (axis cs:28.636,200);
\draw[gray,densely dotted] (axis cs:30,-15) -- (axis cs:30,200);
\node[anchor=north east,font=\scriptsize,align=right] at (axis cs:24.7,196) {metastable band\\ $x_*\le x\le27$};
\end{axis}
\end{tikzpicture}
\caption{Radiatively generated masses of the angular modes of the kernel field at the Standard-Model vacuum (\cref{thm:fullsecond}, \cref{tab:271}). The coloured lines
are the four lowest; they cross zero at $x=27$, $28.5$, $315/11$ and $30$ (dotted). The Standard-Model orbit is a non-degenerate local minimum for $x<27$, and numerically
it is the global minimum for $x<x_*\approx24.884$; in the shaded band it is metastable.}
\label{fig:radiative}
\end{figure}

\begin{proposition}[Random-start minimisation]
\label{prop:randomstart}
$E_x$ was minimised on the unit sphere of $\Kc$ by a quasi-Newton method from Gaussian random starts, with two independent implementations. For $x=0,20,22$ every run
ended at the energy $2x-99$ of the Standard-Model orbit or of its mirror. At $x=25$ a third of the runs ended below it, at a configuration with weight $0.968$ on
$(\mathbf{10},\mathbf1,\mathbf3)$ and $0.032$ on $2(\mathbf6,\mathbf1,\mathbf1)$, whose stabiliser of dimension nine contains $\su2_L$ and projects onto colour with rank six, whereas any
algebra containing $\su3$ projects with rank at least eight: colour is broken. There four right-handed Weyl fermions carry Majorana masses, one of them about twice as heavy
as the other three. Followed in $x$, this minimum crosses the energy of the Standard-Model orbit at $x_*=24.88405$.
\end{proposition}

We conjecture that in the radiative scenario on $\Kc$ the Standard-Model orbit and its mirror are the global minima exactly for $x<x_*$, and that for $x_*<x<27$ the
Standard-Model vacuum is metastable. The competitor again spreads the Majorana mass, now over four right-handed fermions instead of one, which is what the fermion loop
rewards, at the price of breaking colour, which the gauge loop penalises. In terms of masses the local condition and the conjectured global one read
\[
\begin{aligned}
&M_N<27^{1/4}M_{W_R}=2.280\,M_{W_R}\quad\text{and}\quad M_N<2.233\,M_{W_R}&&\text{(one family dominating)},\\
&M_N<9^{1/4}M_{W_R}=1.732\,M_{W_R}\quad\text{and}\quad M_N<1.697\,M_{W_R}&&\text{(three degenerate families)}.
\end{aligned}
\]
They hold under the following restrictions: a tree potential degenerate on $\Kc$ at the matching scale, one loop at leading logarithm, $\kappa=0$, which fermion number
makes automatic, a common gauge coupling for $\su4$ and $\su2_R$, and a single Yukawa coupling for all components of the kernel field (\cref{ass:field}). The masses of
\cref{tab:271} are curvatures of the leading-logarithm effective potential divided by the kinetic norm; pole masses would need the wave-function renormalisation and a
gauge-consistent perturbative expansion, which are not attempted. The bounds are therefore statements about this model and approximation, not bounds on heavy-neutrino
masses outside it. Within it, the result is that gauge fluctuations make the Standard-Model orientation of the
Majorana condensate stable in every direction of the kernel space as long as the heavy neutrinos are not much heavier than $W_R$; beyond that, the local instability is
the condensation of the left-handed triplet, and the configuration that overtakes the Standard-Model orbit globally, at the slightly smaller $x_*$, breaks colour.
